% GENERATED by scripts/build-lecture-slides.py; do not edit. Portable profile: XeTeX, theme and fonts from the core.
\makeatletter\IfFileExists{core/beamer/beamerthememx.sty}{\def\mxcorepath{core/}\def\input@path{{core/beamer/}{core/}}}{\def\mxcorepath{../core/}\def\input@path{{../core/beamer/}{../core/}}}\makeatother
\documentclass[aspectratio=169,10pt,t]{beamer}
\usetheme{mx}
\usepackage{graphicx,booktabs}
\usepackage[french]{babel}
\setlength{\parskip}{.45em}
\hypersetup{pdfauthor={},pdfsubject={Lecture slides},pdfkeywords={},pdftitle={Équations différentielles}}
\providecommand{\N}{\mathbb{N}}
\providecommand{\Z}{\mathbb{Z}}
\providecommand{\D}{\mathbb{D}}
\providecommand{\Q}{\mathbb{Q}}
\providecommand{\R}{\mathbb{R}}
\providecommand{\C}{\mathbb{C}}
\providecommand{\vect}[1]{\overrightarrow{#1}}
\providecommand{\vi}{\vec{\imath}}
\providecommand{\vj}{\vec{\jmath}}
\providecommand{\vk}{\vec{k}}
\providecommand{\norm}[1]{\left\lVert #1\right\rVert}
\providecommand{\vc}[2]{\begin{pmatrix}#1\\#2\end{pmatrix}}
\providecommand{\deter}[2]{\det\left(#1,#2\right)}
\providecommand{\abs}[1]{\left\lvert #1\right\rvert}
\providecommand{\dd}{\,\mathrm{d}}
\providecommand{\e}{\mathrm{e}}
\providecommand{\Card}{\operatorname{Card}}
\providecommand{\E}{\operatorname{E}}
\providecommand{\V}{\operatorname{V}}
\providecommand{\degre}{^{\circ}}
\providecommand{\un}[1]{\,\mathrm{#1}}
\providecommand{\sci}[2]{#1\times10^{#2}}
\providecommand{\pgcd}{\operatorname{PGCD}}
\providecommand{\ppcm}{\operatorname{PPCM}}
\title{Équations différentielles}
\subtitle{Cours et exemples}
\date{}
\author{}
\begin{document}
\mxtitleframe
\begin{frame}{Introduction}
La fonction exponentielle qu'on a étudiée auparavant possède une propriété intéressante :
	      elle est égale à sa propre dérivée.
      Cette propriété va bien souvent être utilisée pour résoudre ce que l'on appelle
      \emph{équation différentielle}.
\end{frame}
\begin{frame}{Introduction}
Une \emph{équation différentielle} est une équation dans laquelle
	      l'inconnue est une fonction et où figurent cette fonction et ses dérivées.
	      Résoudre une équation différentielle sur un intervalle \(I\) consiste donc à chercher toutes les fonctions \(f\)
      (dérivables sur \(I\) autant de fois que nécessaire) qui vérifient l'équation pour tout \(x\) de \(I\).
\end{frame}
\begin{frame}{Équation différentielle du premier ordre}
Une équation différentielle est dite \emph{du premier ordre}
            si elle ne fait intervenir que des dérivées d'ordre au plus \(1\).
            On veut trouver une fonction \(y\) vérifiant l'équation
            \( y' + ay = 0 \) où \(a\) est un nombre réel.
            Autrement dit, on cherche les fonctions \(f\), dérivables sur \(\mathbb{R}\), telles que \({f'(x) + a \; f(x) = 0}\) pour tout réel \(x\).
\end{frame}
\begin{frame}{Équation différentielle du premier ordre}
\textbf{Résolution}

 Cette équation est équivalente à \(y' = - a y\) où encore, en posant \(\alpha = -a\), à
                  \[
                  
                  y' = \alpha y \; .
                  \]
                  La fonction \(f\) définie par \(f(x) = e^{\alpha x}\) vérifie cette propriété
                  puisque \(f'(x) = \alpha \; e^{\alpha x} = \alpha \; f(x)\).

Cette fonction \(f\) est donc une solution particulière de l'Équation eq: diff premier ordre.
\end{frame}
\begin{frame}{Équation différentielle du premier ordre}
Cherchons maintenant \emph{toutes} les solutions.
                  Soit \(f\) une solution quelconque de eq: diff premier ordre.
	                  Comme \(e^{\alpha x} \neq 0\) pour tout réel \(x\), on peut définir sur \(\mathbb{R}\) la fonction \(g\) par
	                  \(g(x) = f(x) \, e^{-\alpha x}\), de sorte que \(f(x) = g(x) \, e^{\alpha x}\) :
	                  toute fonction peut donc s'écrire sous cette forme, ce n'est pas une hypothèse.
	                  La fonction \(g\) est dérivable comme produit de fonctions dérivables, et
                  \[\begin{aligned}
                        g'(x) & = f'(x) e^{-\alpha x } - \alpha f(x) e^{-\alpha x}
                        \\ \text{et comme $f'(x) = \alpha f(x)$,} \\
                        g'(x) & = \alpha f(x) e^{-\alpha x } - \alpha f(x) e^{-\alpha x} = 0 \;.
                  \end{aligned}\]
	                  La dérivée de \(g\) est nulle sur l'intervalle \(\mathbb{R}\), donc \(g\) est constante :
	                  il existe un réel \(k\) tel que \(g(x) = k\), c'est-à-dire \(f(x) = k \, e^{\alpha x}\) pour tout \(x\).
	                  Réciproquement, pour tout réel \(k\), la fonction \(x \mapsto k\,e^{\alpha x}\) est solution,
	                  car sa dérivée est \(k \alpha e^{\alpha x}\).
	                  Ainsi :
\end{frame}
\begin{frame}{Équation différentielle du premier ordre}
l'ensemble des solutions de l'équation \(y' + a \; y = 0\) est
	                  l'ensemble des fonctions \(f\) telles que \(f(x) = k \; e^{-a x}\) où \(k \in \mathbb{R}\).
\end{frame}
\begin{frame}{Équation différentielle du premier ordre}
\textbf{Résolution d'une équation avec condition initiale}

 On veut résoudre une équation différentielle du premier ordre
                  vérifiant une condition de la forme \(y(x_0) = y_0\).
                  On sait déjà que les solutions de l'équation sont de la forme \(y(x) = k \; e^{-a x}\).
                  Ainsi, \(y_0 = y(x_0) = k \; e^{-a x_0}\), ce qui nous donne \(k = y_0 e^{a x_0}\).
                  La fonction \(y\) s'écrit alors
                  \(y(x) = y_0 e^{a x_0} e^{-a x} = y_0 e^{-a(x - x_0)}\).
\end{frame}
\begin{frame}{Équation différentielle du premier ordre (2)}
Le calcul précédent montre que la constante \(k\) est entièrement déterminée :
                  l'équation \(y' + a y = 0\) admet une
                  solution unique vérifiant la condition \(y(x_0) = y_0\).
\end{frame}
\begin{frame}{Équation différentielle du premier ordre}
\textbf{Équation avec second membre}

 Soit \(h\) une fonction définie sur un intervalle \(I\) et l'équation
                  \((E) : y' + a y = h(x)\). L'équation \((E_0) : y' + a y = 0\) est appelée
                  \emph{équation sans second membre} (ou homogène) associée à \((E)\).
\end{frame}
\begin{frame}{Équation différentielle du premier ordre (2)}
Supposons que l'on connaisse une solution particulière \(p\) de \((E)\).
                  Pour une fonction \(f\) dérivable sur \(I\), posons \(u = f - p\). On a
                  \(u' + a u = (f' + a f) - (p' + a p) = (f' + a f) - h\).
                  Donc \(f\) est solution de \((E)\) si et seulement si \(u\) est solution de \((E_0)\),
                  c'est-à-dire si et seulement si \(u(x) = k e^{-ax}\) :

les solutions de \(y' + ay = h(x)\) sont les fonctions \(x \mapsto p(x) + k\,e^{-ax}\), \(k \in \mathbb{R}\).
\end{frame}
\begin{frame}{Équation différentielle du premier ordre}
En particulier, si \(a \neq 0\) et si \(h\) est une constante \(b\), la fonction constante \(p = \dfrac{b}{a}\)
                  est une solution particulière, et les solutions de \(y' + a y = b\) sont les fonctions
                  \(x \mapsto \dfrac{b}{a} + k\,e^{-ax}\).
                  Le même raisonnement s'applique aux équations du second ordre de la section suivante :
                  solution générale = solution particulière + solution générale de l'équation sans second membre.
\end{frame}
\begin{frame}{Équation différentielle du second ordre}
On s'intéresse maintenant aux équations de la forme
            \[
            
                  y'' + a y' + by = 0,
            \]
            où \(a\) et \(b\) sont des nombres réels.
            Pour résoudre une telle équation, on résout d'abord l'équation caractéristique.
\end{frame}
\begin{frame}{Équation différentielle du second ordre}
\textbf{Équation caractéristique:}

 On cherche les réels \(r\) tels que la fonction \(x \mapsto e^{rx}\) soit une solution de l'équation.
                  Posons \(f(x) = e^{r x}\), on a alors \(f'(x) = r e^{rx}\) et \(f''(x) = r^2 e^{rx}\).
                  Ainsi, \(f''(x) + a f'(x) + b f(x) = 0\) équivaut à
                  \(r^2 e^{rx} + a r e^{rx} + b e^{r x} = 0\) et par conséquent à \((r^2 + a r + b) e^{r x} = 0\).
	                  Comme la fonction exponentielle est toujours strictement positive (et donc non nulle),
                  cela équivaut à \(r^2 + a r + b = 0\).
                  Cette équation du second degré est appelée \emph{équation caractéristique} de eq: diff second ordre.
\end{frame}
\begin{frame}{Équation différentielle du second ordre}
\textbf{Résolution:}

 Selon le discriminant \(\Delta = a^2 - 4b\) de l'équation caractéristique, on admet le théorème suivant.
\end{frame}
\begin{frame}{Équation différentielle du second ordre}
\begin{itemize}
\item Si \(\Delta > 0\):
                                    Soient \(r_1\) et \(r_2\) les deux solutions réelles de l'équation caractéristique.
                                    L'ensemble des solutions de l'Équation eq: diff second ordre est
                                    l'ensemble des fonctions \(f\) de la forme
                                    \(f(x) = A \; e^{r_1 \; x } + B \; e^{r_2 \; x}\), où \(A\) et \(B\) sont des réels.

\end{itemize}
\end{frame}
\begin{frame}{Équation différentielle du second ordre (2)}
\begin{itemize}
\item Si \(\Delta = 0\): Soit \(\rho\) la racine (double) de l'équation caractéristique.
                                    L'ensemble des solutions de l'Équation eq: diff second ordre est
                                    l'ensemble des fonctions \(f\) de la forme
	                                    \(f(x) = (A x + B) \; e^{\rho x} \), où \(A\) et \(B\) sont des réels.

\end{itemize}
\end{frame}
\begin{frame}{Équation différentielle du second ordre (3)}
\begin{itemize}
\item Si \(\Delta < 0\):
                                    L'équation caractéristique admet alors deux racines complexes conjuguées.
                                    Soient \(\alpha\) et \(\beta\) les deux réels (\(\beta \neq 0\)) tels que
                                    les deux solutions sont \(\alpha + i \beta\) et \(\alpha - i \beta\).
                                    L'ensemble des solutions de l'Équation eq: diff second ordre est
                                    l'ensemble des fonctions \(f\) de la forme
                                    \(f(x) = \left(A \cos (\beta x ) + B \sin (\beta x) \right) \; e^{\alpha x}\), où \(A\) et \(B\) sont des réels.

\end{itemize}
\end{frame}
\begin{frame}{Équation différentielle du second ordre}
Dans tous les cas, il existe une solution unique satisfaisant aux conditions initiales
                        \(y(x_0) = y_0\) et \(y'(x_0) = y'_0\).
\end{frame}
\begin{frame}{Équation différentielle du second ordre}
Résoudre l'équation différentielle
                        \[
                        
                              y'' - 3 y' - 4 y = 0
                        \]
                  avec les conditions initiales : \(y(0) = 2 \) et \(y'(0) = 4\).
\end{frame}
\begin{frame}{Équation différentielle du second ordre (2)}
L'équation caractéristique est \(r^2 - 3r - 4 = 0\) ; son discriminant est \(\Delta = 9 + 16 = 25 > 0\)
                  et ses racines sont \(r_1 = \dfrac{3 + 5}{2} = 4\) et \(r_2 = \dfrac{3 - 5}{2} = -1\).
                  Les solutions de eq:ex:eq diff 1 sont donc les fonctions \(y(x) = A e^{4x} + B e^{-x}\),
                  de dérivée \(y'(x) = 4A e^{4x} - B e^{-x}\).
                  Les conditions initiales donnent le système \(A + B = 2\) et \(4A - B = 4\),
                  d'où \(5A = 6\), \(A = \dfrac{6}{5}\) et \(B = \dfrac{4}{5}\).
                  La solution cherchée est \(y(x) = \dfrac{6}{5} e^{4x} + \dfrac{4}{5} e^{-x}\).
\end{frame}
\end{document}
